\begin{textsample}{POS Dim 7 – global\_north\_2024\_11 – Score 2.00 – global\_north\_2024\_11/117891420.txt}  \label{ex:f7_pos_196}
BEIRUT : Hezbollah ' s leader delivered a defiant speech on Wednesday ( Nov 20 ) saying Israel can not impose its conditions for a truce, as US envoy Amos Hochstein headed from Lebanon to Israel to try to end the war.
Israeli Foreign Minister Gideon Saar, in a near-simultaneous statement, said any ceasefire deal must ensure Israel has the"freedom to act"against Hezbollah.
Hochstein announced in Lebanon that he would head to Israel on Wednesday to try to seal a ceasefire agreement in the war in Lebanon, which escalated in late September after nearly a year of deadly exchanges of fire across Israel ' s northern border.
Israel expanded the focus of its operations from Gaza to Lebanon, vowing to secure the north and allow tens of thousands of people displaced by the cross-border fire to return home.
It has also intensified strikes on neighbouring Syria, a key \textbf{conduit} of weapons for Hezbollah from its backer Iran.
A woman rests in an ambulance at the site where a rocket, fired @ @ @ @ @ @ @ @ @ @ Israel, Wednesday, Nov 20, 2024. ( Photo : AP/Francisco Seco ) In the latest reported attack, the Syrian defence ministry said 36 people were killed and more than 50 wounded in israeli strikes on the oasis city of Palmyra."Israel can not defeat us and can not impose its conditions on us,"Hezbollah leader Naim Qassem said in an address broadcast shortly after Hochstein announced he would travel to Israel.
Qassem added that his armed group seeks a"complete and comprehensive end to the aggression"and"the preservation of Lebanon ' s sovereignty".
He also vowed that the response to recent deadly Israeli strikes on Beirut would be on"central Tel Aviv", Israel ' s densely populated commercial hub.
Related :"PROGRESS"Before heading to Israel, Hochstein met for a second time with one of his main interlocutors, Hezbollah-allied parliament speaker Nabih Berri, who has led mediation efforts on behalf of the Iran-backed group."The meeting today built @ @ @ @ @ @ @ @ @ @ I will travel from here in a couple hours to Israel to try to bring this to a close if we can,"Hochstein told reporters in the Lebanese capital.
Hochstein had on Tuesday said an end to the war was"within our grasp", while a diplomat in Lebanon told AFP that he had studied some modifications to the US truce plan with Lebanese officials.
Ahead of Hochstein ' s arrival, Israel ' s top diplomat Saar said :"In any agreement we will reach, we will need to keep the freedom to act if there will be violations."Related : Among them were more than 200 children, according to the United Nations.
While Hochstein was in Beirut, the situation in the capital was relatively calm Tuesday and Wednesday, but south Lebanon, where Hezbollah holds sway, has seen battles and strikes.
The United States, Israel ' s main military and political backer, has been pushing for a UN \textbf{resolution} that ended the last Hezbollah-Israel war in 2006 @ @ @ @ @ @ @ @ @ @ Under UN Security Council \textbf{Resolution} 1701, Lebanese troops and UN peacekeepers should be the only armed forces deployed in south Lebanon.
Lebanese Army Intelligence members inspect an army position that was damaged in an Israeli airstrike on Tuesday night, in Sarafand, southern Lebanon, Wednesday, Nov 20, 2024. ( Photo : AP/Mohammed Zaatari ) VIOLENCE IN SOUTH LEBANON While not engaged in the ongoing war, the Lebanese army has reported 18 fatalities from among its ranks since Sep 23.
On Wednesday, the army said Israeli fire killed a soldier in south Lebanon, a day after it announced the deaths of three other personnel in a strike.
The Israeli military later said, without mentioning the deaths, that it was looking into reports of Lebanese soldiers injured by a strike on Tuesday."We emphasise that the ( Israeli army ) is operating precisely against the Hezbollah terrorist organisation and is not operating against the Lebanon Armed Forces,"the military told AFP in a statement.
Lebanon ' s @ @ @ @ @ @ @ @ @ @ in south Lebanon overnight and on Wednesday, saying Israeli troops were seeking to advance further near the town of Khiam.
Hezbollah said Wednesday that it had twice targeted Israeli troops near the flashpoint border town, home to an infamous former detention centre that was shut down after the end of the Israeli occupation of south Lebanon in 2000.
The NNA said that Israeli forces were"attempting to advance from the Kfarshuba hills... to open up a new front under the cover of fire and artillery shells and air strikes"."Violent clashes are taking place"between Hezbollah and Israeli forces, it added.
Israel said Wednesday it hit 100"terror targets"around Lebanon in the past day, including"launchers, weapons storage facilities, command centres and military structures".
Hezbollah, meanwhile, said it had launched drones at two Israeli military bases in northern Israel and fired rockets at the town of Safed.
Source : AFP/fs Sign up for our newsletters Get our pick @ @ @ @ @ @ @ @ @ @

% matched lemmas: conduit, resolution
\end{textsample}
